\section{Height scales and an auxiliary sequence theorem}
\label{sec:sequences}

The proof separates coordinates whose heights grow at different rates.
We first explain the linear algebra that permits real homomorphisms to act
on normalized points. We then state the sequence theorem and show that it
implies Theorem~\ref{thm:nondensity}. Its proof occupies the subsequent
sections.

\subsection{Real homomorphisms and height norms}

For an abelian variety $G$ over $\Qbar$, choose a symmetric ample line
bundle and write
\[
 \mathcal H(G)=G(\Qbar)\otimes_{\mathbb Z}\mathbb R,
 \qquad \|x\|^2=\widehat h(x).
\]
The canonical height pairing extends to a positive-definite real
quadratic form on $\mathcal H(G)$. Indeed, every finite collection of
points is defined over one number field, and the assertion on its real
span is the positive definiteness of the N\'eron--Tate pairing on the
Mordell--Weil group modulo torsion. We use the same notation for a point
and its image in $\mathcal H(G)$. In particular, torsion points have image
zero. Polarizations also give Hermitian forms on the complex Lie algebras
$\Lie(G)$. On products we choose product polarizations and product norms.
We use the standard complex uniformization, complete reducibility,
finite-rank homomorphism groups and polarization adjoints recalled in
\cite[Chapter~I, Sections~2, 10 and 14]{Milne2008AV}. Canonical metrics
and heights are taken with the multiplication normalization of
\cite[Theorem~2.2]{Zhang1995}.

\begin{lemma}[Real homomorphism calculus]
\label{lemma:real-calculus}
Let $G,G'$ be abelian varieties over $\Qbar$ with the preceding choices.
Every element of
$\Hom(G,G')_{\mathbb R}=\Hom(G,G')\otimes_{\mathbb Z}\mathbb R$
acts on both $\Lie(G)$ and $\mathcal H(G)$. The Lie action is faithful.
Convergence in this finite-dimensional real homomorphism space implies
operator-norm convergence on the height spaces.

The polarization adjoint of a real homomorphism is a real homomorphism.
Its kernel and image in the Lie algebras have orthogonal projectors in
the corresponding real endomorphism algebras, and it has a generalized
inverse in the opposite real homomorphism space. In particular, if $f$
is fixed and $\|f(x_i)\|\to0$, then
\[
 \|P_{(\ker f)^\perp}(x_i)\|\longrightarrow0.
\]
\end{lemma}

\begin{proof}
Fix integral homomorphisms $f_1,\ldots,f_b$ forming a real basis of
$\Hom(G,G')_{\mathbb R}$. For each $j$, ample domination of $f_j^*L'$
by a power of $L$ gives
$\widehat h_{L'}(f_j(x))\le C_j\widehat h_L(x)$ for all algebraic
points $x$. The same inequality holds on $\mathcal H(G)$, by its
quadratic-form interpretation on every finite-dimensional span. Hence
\[
 \left\|\sum_j t_j f_j(x)\right\|
 \le \left(\sum_j |t_j|\sqrt{C_j}\right)\|x\|.
\]
The constants are independent of the fields of definition of the
points. This proves the assertion about operator norms.

The underlying real Lie representation identifies with the real
extension of the representation on rational homology. It is faithful,
also after scalar extension from $\mathbb Q$ to $\mathbb R$. The
polarization adjoint belongs to rational homomorphisms, and therefore
extends to real homomorphisms. For a real homomorphism $f$, the operator
$f^*f$ is nonnegative and self-adjoint on $\Lie(G)$. Choose a real
polynomial $p$ that is zero at zero and equals $1/t$ at every positive
eigenvalue $t$ of $f^*f$. Then
\[
 f^\dagger=p(f^*f)f^*
\]
is a real homomorphism and
$f^\dagger f=P_{(\ker f)^\perp}$,
$ff^\dagger=P_{\operatorname{im}f}$.
Faithfulness makes these identities identities of real homomorphisms,
so they also hold on the height spaces. The last assertion follows by
applying the fixed bounded operator $f^\dagger$.
\end{proof}

In particular, the image of any real homomorphism into $G$ is the image
of a real idempotent endomorphism of $G$. Images, inverse images, and
intersections of finitely many such Lie subspaces can therefore be
handled by real homomorphisms. These subspaces need not be Lie algebras
of abelian subvarieties.

\subsection{The sequence theorem}

Local fiber dimension at a point always means the maximum of the
dimensions of the irreducible components of the fiber through that
point. This convention matters when a marked point is singular.

The fiber bounds in the sequence theorem record the dimension
contributed by each height level. Let $W_1,W_2$ be abelian varieties.
Consider an irreducible $X\subseteq W_1\times W_2$, points
$(p_{i1},p_{i2})\in X$, and scales $H_{i\ell}\ge1$ with
$\|p_{i\ell}\|=O(\sqrt{H_{i\ell}})$ and
$H_{i1}/H_{i2}\to0$. Retain the first coordinate when forming the
second ambient factor:
\[
 \begin{array}{lll}
 A_1=W_1,&Y_1=\operatorname{pr}_1(X),&m_1=\dim Y_1,\\
 A_2=W_1\times W_2,&Y_2=X,&m_2=\dim X-\dim Y_1.
 \end{array}
\]
The distinguished subgroups are $W_1\subseteq A_1$ and
$\{0\}\times W_2\subseteq A_2$. The integers $m_1,m_2$ are the
generic dimensions of the corresponding fibers, and
$m_1+m_2=\dim X$. The copied coordinate
$p_{i1}/\sqrt{H_{i2}}$ tends to zero, so it disappears after
normalization at the second scale. In the reduction below, genericity
of the sequence ensures that the marked points eventually have these
fiber dimensions. The following definition allows arbitrary pairs
$W_\ell\subseteq A_\ell$, accommodating both this construction with
any number of levels and the quotients used in Section~\ref{sec:tangent}.

\begin{definition}[Sequence data]
\label{def:sequence-data}
Let $I$ be the ordered list $1,\ldots,s$, possibly preceded by $0$.
The list is nonempty; if $0$ occurs, $s$ may be zero. For each
$\ell\in I$, fix an abelian variety $A_\ell$ over $\Qbar$, a nonzero
abelian subvariety $W_\ell\subseteq A_\ell$, an irreducible subvariety
$Y_\ell\subseteq A_\ell$, and an integer $m_\ell\ge0$. Put
\[
 \cA=\prod_{\ell\in I}A_\ell,
 \qquad W=\prod_{\ell\in I}W_\ell,
 \qquad D_\ell=A_\ell/W_\ell.
\]
For positive integers $i$, let $y_{i\ell}\in Y_\ell(\Qbar)$ and
$H_{i\ell}\ge1$ satisfy
\begin{equation}
 \label{eq:sequence-fibers}
 \dim_{y_{i\ell}}\bigl(Y_\ell\cap(y_{i\ell}+W_\ell)\bigr)
 \le m_\ell,
 \qquad \|y_{i\ell}\|=O(\sqrt{H_{i\ell}}).
\end{equation}
For consecutive indices $\ell,\ell'$ in $I$, require
$H_{i\ell}/H_{i\ell'}\to0$.

The factors $\ell\ge1$ are called high factors. They satisfy
$H_{i\ell}\to\infty$, and, for every fixed homomorphism
$g:A_\ell\to G$ to an abelian variety whose restriction to $W_\ell$
is nonzero,
\begin{equation}
 \label{eq:sequence-nondegenerate}
 \liminf_{i\to\infty}
 \frac{\|g(y_{i\ell})\|}{\sqrt{H_{i\ell}}}>0.
\end{equation}
If the index $0$ occurs, it is called the low factor and satisfies
$W_0=A_0$, $H_{i0}=1$. Its points eventually avoid each fixed proper
torsion coset of $A_0$.
\end{definition}

The homomorphism in \eqref{eq:sequence-nondegenerate} is integral; the
condition is equivalently unchanged by allowing rational homomorphisms.
Its positive lower bound may depend on $g$. No bound on the degrees of
the marked points is assumed.

\begin{theorem}[Auxiliary sequence theorem]
\label{thm:sequence}
For sequence data as in Definition~\ref{def:sequence-data}, let
$V\subseteq\Lie(W)$ be the image of an idempotent in
$\End(\cA)_{\mathbb R}$. Let $\Phi$ be the orthogonal projector with
kernel $V$, and put
\[
 q_i=(y_{i\ell}/\sqrt{H_{i\ell}})_{\ell\in I}
 \in\mathcal H(\cA).
\]
The following two conditions cannot both hold:
\begin{equation}
 \label{eq:forbidden-sequence}
 \codim_{\Lie(W)}V>\sum_{\ell\in I}m_\ell,
 \qquad \|\Phi(q_i)\|\longrightarrow0.
\end{equation}
\end{theorem}

We will first prove this theorem for all-high data in every dimension.
We then treat data with a low factor. At each stage, a hypothetical
counterexample of least $\dim\cA$ gives the geometric exclusion proved
in Section~\ref{sec:tangent}. The order of these two stages permits a
projection that removes the low factor.

\subsection{Decomposing a sequence into height scales}

\begin{lemma}[Height decomposition]
\label{lemma:height-decomposition}
Let $A$ be a nonzero abelian variety and $x_i\in A(\Qbar)$. After
passing to a subsequence, there is an isogeny
\[
 \alpha:A\longrightarrow\prod_{\ell\in I}W_\ell
\]
with nonzero factors, and scales $H_{i\ell}\ge1$, such that the
coordinates $p_{i\ell}$ of $\alpha(x_i)$ have norms
$O(\sqrt{H_{i\ell}})$ and successive scale ratios tend to zero.
There is at most one bounded scale, placed first and equal to one.
Every other scale tends to infinity and satisfies
\eqref{eq:sequence-nondegenerate} with $A_\ell=W_\ell$ for every
fixed nonzero homomorphism on $W_\ell$.
\end{lemma}

\begin{proof}
If the point heights are bounded, use $A$ itself as the bounded factor.
Otherwise take a subsequence with $H_i=\|x_i\|^2\to\infty$.
Among fixed quotient homomorphisms $\pi:A\to Q$ for which
$\|\pi(x_i)\|^2=o(H_i)$ on some subsequence, choose one with
maximal $\dim Q$, and restrict to that subsequence. The zero quotient
is allowed. A quotient of dimension $\dim A$ is an isogeny and cannot
have this property, by height comparison. Poincar\'e reducibility
provides an isogeny $A\to Q\times W'$ whose first component is
$\pi$; harmless fixed isogenies of the targets can be absorbed here.

Write $p'_i$ for the $W'$ coordinate. Suppose a fixed nonzero
homomorphism $g$ on $W'$ satisfies
\[
 \liminf_{i\to\infty}\frac{\|g(p'_i)\|}{\sqrt{H_i}}=0.
\]
Take a further subsequence on which this ratio tends to zero. The homomorphism
whose coordinates are $\pi$ and $g$ on the complementary factor has
image of dimension $\dim Q+\dim g(W')>\dim Q$, and its point
heights are $o(H_i)$. Quotienting by the identity component of its
kernel changes its image only by an isogeny. This contradicts the
maximality of $Q$. Thus $W'$ is a high factor with scale $H_i$ and the
required nondegeneracy for every fixed homomorphism.

If $Q=0$, stop. Otherwise continue on $Q$ with the points
$\pi(x_i)$. If their heights are
bounded, this is the low factor; if not, extract a diverging scale and
repeat. At each step the remaining quotient has smaller dimension and
its squared heights are little-oh of the preceding scale. The process
therefore terminates. Reverse the order of the extracted factors, so
that the scales are increasing. Composing the finitely many fixed
isogenies gives $\alpha$. Fixed height comparisons preserve all the
asserted bounds and ratios.
\end{proof}

\begin{proposition}
\label{prop:sequence-implies-nondensity}
The auxiliary sequence theorem implies Theorem~\ref{thm:nondensity}
for every abelian variety over $\Qbar$.
\end{proposition}

\begin{proof}
Let $X\subseteq A$ be irreducible and contained in no proper torsion
coset. If $\dim A=0$, there are no torsion cosets of codimension
$\dim X+1$, so there is nothing to prove. Suppose otherwise that
$X\cap A^{[\dim X+1]}$ is Zariski dense. Since $\Qbar$ is
countable, choose a sequence in this intersection that eventually
avoids every fixed proper closed subset of $X$ defined over $\Qbar$.
Apply Lemma~\ref{lemma:height-decomposition} and replace $A,X$ and
the sequence by their isogeny images. Dimensions of subvarieties and
codimensions of torsion cosets are preserved by this finite map. Its
image $X$ is still contained in no proper torsion coset. We have now
written $A=\prod_{\ell\in I}W_\ell$ and
$x_i=(p_{i\ell})_\ell$.

For each $i$, choose a torsion coset $B_i+\tau_i$ containing $x_i$
with $\codim_A B_i>\dim X$. After extraction its codimension is
constant. In $\Lie(A)$ let
\[
 U_i=\operatorname{diag}(H_{i\ell}^{-1/2})\Lie(B_i).
\]
The orthogonal projector $P_i$ onto $U_i$ is a real endomorphism by
Lemma~\ref{lemma:real-calculus}, applied to the scaled subgroup
inclusion. All these projectors have the same rank. A subsequence
converges to an orthogonal projector $P$ of the same rank. Moreover
$P\in\End(A)_{\mathbb R}$, because this finite-dimensional
subspace is closed in the space of real linear maps on $\Lie(A)$.
Put $U=\operatorname{im}P$. Then
\begin{equation}
 \label{eq:initial-limit-projector}
 \codim_{\Lie(A)}U>\dim X,
 \qquad
 \left\|(1-P)(p_{i\ell}/\sqrt{H_{i\ell}})_\ell\right\|
 \longrightarrow0.
\end{equation}
For the second assertion, torsion vanishes in the height space, so
$1-P_i$ kills the indicated normalized vector exactly. Those vectors
are bounded, and Lemma~\ref{lemma:real-calculus} converts convergence
of the projectors into degree-independent operator-norm convergence.

To obtain the local fiber bounds in the sequence theorem, we retain
copies of all slower coordinates. For each $\ell\in I$ set
\[
 A_\ell=\prod_{h\le\ell}W_h,
 \qquad Y_\ell=\operatorname{pr}_{\le\ell}(X),
 \qquad y_{i\ell}=(p_{ih})_{h\le\ell}.
\]
Regard $W_\ell$ as the last coordinate in $A_\ell$. Let
$m_\ell=\dim Y_\ell-\dim Y_{\ell^-}$, where $Y_{\ell^-}$ is
the preceding projection and is a point for the first index. These
integers are nonnegative and sum to $\dim X$. The local dimension of
the fiber of $Y_\ell\to Y_{\ell^-}$ equals $m_\ell$ on a
nonempty open subset. The sequence eventually lies in that subset:
the inverse image in $X$ of its complement is a fixed proper closed
set. Thus \eqref{eq:sequence-fibers} holds after discarding finitely
many terms.

For a high factor, every slower coordinate is
$o(\sqrt{H_{i\ell}})$ in norm. A fixed homomorphism on $A_\ell$
nonzero on its last $W_\ell$ coordinate is the sum of a nonzero
homomorphism on that coordinate and fixed homomorphisms on the slower
ones. The triangle inequality and
Lemma~\ref{lemma:height-decomposition} prove
\eqref{eq:sequence-nondegenerate}. If a low factor occurs, its image
of $X$ is contained in no proper torsion coset: the inverse image of
such a coset is a finite union of proper torsion cosets, one of which
would contain irreducible $X$. The same pullback argument proves
eventual avoidance for the low sequence.

Finally let $j:A\to\cA$ insert each $W_\ell$ into the distinguished
last coordinate of $A_\ell$, and let $\rho:\cA\to A$ be the
coordinate retraction. Put $V=j(U)$. Since $\rho j=1$, the real
endomorphism $jP\rho$ is idempotent with image $V$, and
\[
 \codim_{\Lie(W)}V=\codim_{\Lie(A)}U>\dim X
 =\sum_\ell m_\ell.
\]
Write $q_i^0=(p_{i\ell}/\sqrt{H_{i\ell}})_\ell$ for the normalized
vector before enlargement, and $q_i$ for the enlarged vector in
Theorem~\ref{thm:sequence}. Their difference has only copied slower
coordinates. With product norms,
\[
 \|q_i-jq_i^0\|^2
 =\sum_\ell\sum_{h<\ell}\frac{\|p_{ih}\|^2}{H_{i\ell}}
 =O\!\left(\sum_\ell\sum_{h<\ell}
                  \frac{H_{ih}}{H_{i\ell}}\right)\longrightarrow0.
\]
If $\Phi$ is the complementary projector for $V$, then
$\Phi j=\Phi j(1-P)$. Equation~\eqref{eq:initial-limit-projector}
and the displayed error bound therefore give $\|\Phi q_i\|\to0$.
We have constructed \eqref{eq:forbidden-sequence}, a contradiction to
Theorem~\ref{thm:sequence}.
\end{proof}
